\subtitle{\bf Aula 13 - A Linguagem Line}
\begin{frame}[t,plain]
\titlepage
\end{frame}

\begin{frame}{Objetivos}
\protect\hypertarget{objetivos-1}{}
\begin{itemize}
\item
  Apresentar a semântica operacional para uma linguagem imperativa
  simples.
\item
  Apresentar a implementação do interpretador desta linguagem.
\end{itemize}
\end{frame}



\begin{frame}{Motivação: Expressões Puras Não Bastam}
\protect\hypertarget{motivauxe7uxe3o-expressuxf5es-puras-nuxe3o-bastam}{}
\begin{itemize}

\item
  No capítulo anterior: expressões \textbf{puras} sem variáveis
\item
  O valor de uma expressão dependia apenas da sua estrutura sintática
\end{itemize}
\end{frame}

\hypertarget{a-linguagem-line-1}{%
\section{A Linguagem Line}\label{a-linguagem-line-1}}

\begin{frame}{Motivação: Expressões Puras Não Bastam}
\protect\hypertarget{motivauxe7uxe3o-expressuxf5es-puras-nuxe3o-bastam-1}{}
\begin{itemize}

\item
  Programas reais manipulam \textbf{estado}: valores armazenados em
  variáveis
\item
  Precisamos estender a semântica com um componente de estado
\end{itemize}
\end{frame}

\begin{frame}{Sintaxe da Linguagem Line}
\protect\hypertarget{sintaxe-da-linguagem-line}{}
\[\begin{array}{lcl}
\mathit{prog} &::=& s^* \\[4pt]
s &::=& x \mathbin{:=} e \\
  &\mid& \mathbf{read}\; x \\
  &\mid& \mathbf{print}\; e \\[4pt]
e &::=& n \mid x \mid e_1 + e_2 \mid e_1 * e_2
\end{array}\]
\end{frame}

\begin{frame}[fragile]{Tipos em Haskell}
\protect\hypertarget{tipos-em-haskell}{}
\begin{lstlisting}[language=Haskell]
data Line = Line [Stmt]

data Stmt
  = SAssign Var Exp
  | SRead   Var
  | SPrint  Exp
\end{lstlisting}
\end{frame}

\begin{frame}[fragile]{Tipos em Haskell}
\protect\hypertarget{tipos-em-haskell-1}{}
\begin{lstlisting}[language=Haskell]
data Exp
  = EInt Int
  | EVar Var
  | Exp :+: Exp
  | Exp :*: Exp

type Var = String
\end{lstlisting}
\end{frame}

\hypertarget{ambientes}{%
\section{Ambientes}\label{ambientes}}

\begin{frame}[fragile]{Definição de Ambiente}
\protect\hypertarget{definiuxe7uxe3o-de-ambiente}{}
Um \textbf{ambiente}
\(\sigma : \mathit{Var} \rightharpoonup \mathbb{Z}\) é uma função
parcial que mapeia nomes de variáveis para inteiros.

Em Haskell: \texttt{type Env = Map String Int}
\end{frame}

\begin{frame}{Notação}
\protect\hypertarget{notauxe7uxe3o}{}
\begin{itemize}

\item
  \(\sigma(x)\): valor de \(x\) em \(\sigma\) (supondo
  \(x \in \mathrm{dom}(\sigma)\))
\item
  \(\sigma[x \mapsto n]\): ambiente obtido atualizando \(x\) com \(n\)
\item
  \(\emptyset\): ambiente vazio (sem variáveis definidas)
\end{itemize}
\end{frame}

\begin{frame}[fragile]{Por que Precisamos de Ambientes?}
\protect\hypertarget{por-que-precisamos-de-ambientes}{}
\begin{itemize}

\item
  Expressões agora podem conter \textbf{variáveis}:
  \texttt{x * 4}
\item
  O valor de \texttt{x} só é conhecido durante a
  execução
\end{itemize}
\end{frame}

\begin{frame}{Por que Precisamos de Ambientes?}
\protect\hypertarget{por-que-precisamos-de-ambientes-1}{}
\begin{itemize}

\item
  O ambiente \(\sigma\) registra o estado da memória em cada instante
\item
  Comandos de atribuição \textbf{modificam} o ambiente:
  \(\sigma \to \sigma[x \mapsto n]\)
\end{itemize}
\end{frame}

\begin{frame}{Por que Precisamos de Ambientes?}
\protect\hypertarget{por-que-precisamos-de-ambientes-2}{}
\begin{itemize}

\item
  Expressões apenas \textbf{consultam} o ambiente, sem modificá-lo
\end{itemize}
\end{frame}

\hypertarget{semuxe2ntica-small-step}{%
\section{Semântica Small-Step}\label{semuxe2ntica-small-step}}

\begin{frame}{Expressões com Variáveis}
\protect\hypertarget{expressuxf5es-com-variuxe1veis}{}
\begin{itemize}

\item
  A configuração de uma expressão é um par \((\sigma, e)\); a redução é
  \((\sigma, e) \to (\sigma, e')\).
\end{itemize}
\end{frame}

\begin{frame}{Expressões com Variáveis}
\protect\hypertarget{expressuxf5es-com-variuxe1veis-1}{}
\begin{itemize}

\item
  \textbf{Variável:}
\end{itemize}

\[\frac{x \in \mathrm{dom}(\sigma)}{(\sigma,\, x) \to (\sigma,\, \sigma(x))} \quad\text{(Var)}\]
\end{frame}

\begin{frame}{Expressões com Variáveis}
\protect\hypertarget{expressuxf5es-com-variuxe1veis-2}{}
\begin{itemize}

\item
  Uma variável é substituída pelo seu valor corrente.
\item
  Se \(x \notin \mathrm{dom}(\sigma)\): \textbf{erro} de variável não
  inicializada.
\end{itemize}
\end{frame}

\begin{frame}{Regras para Expressões}
\protect\hypertarget{regras-para-expressuxf5es}{}
\[\frac{(\sigma, e_1) \to (\sigma, e_1')}{(\sigma,\, e_1 + e_2) \to (\sigma,\, e_1' + e_2)} \quad\text{(AddL)}
\]
\end{frame}

\begin{frame}{Regras para Expressões}
\protect\hypertarget{regras-para-expressuxf5es-1}{}
\[
\frac{(\sigma, e_2) \to (\sigma, e_2')}{(\sigma,\, n_1 + e_2) \to (\sigma,\, n_1 + e_2')} \quad\text{(AddR)}\]
\end{frame}

\begin{frame}{Regras para Expressões}
\protect\hypertarget{regras-para-expressuxf5es-2}{}
\[\frac{}{(\sigma,\, n_1 + n_2) \to (\sigma,\, n_1 \mathbin{\hat{+}} n_2)} \quad\text{(Add)}\]

\begin{itemize}

\item
  Note que as expressões \textbf{não modificam} \(\sigma\)

  \begin{itemize}

  \item
    São puras no sentido de que apenas consultam o ambiente.
  \end{itemize}
\end{itemize}
\end{frame}

\begin{frame}{Regras para Comandos}
\protect\hypertarget{regras-para-comandos}{}
A configuração de um comando é \((\sigma, s) \to (\sigma', s')\), onde
\(s'\) é o resíduo após a redução.
\end{frame}

\begin{frame}{Regras para Comandos}
\protect\hypertarget{regras-para-comandos-1}{}
\begin{itemize}

\item
  \textbf{Atribuição (avalia primeiro):}
  \[\frac{(\sigma, e) \to (\sigma, e')}{(\sigma,\; x \mathbin{:=} e) \to (\sigma,\; x \mathbin{:=} e')} \quad\text{(AssignStep)}\]
\end{itemize}
\end{frame}

\begin{frame}{Regras para Comandos}
\protect\hypertarget{regras-para-comandos-2}{}
\begin{itemize}

\item
  \textbf{Atribuição (valor final):}
  \[\frac{}{(\sigma,\; x \mathbin{:=} n) \to (\sigma[x \mapsto n],\; \mathbf{skip})} \quad\text{(Assign)}\]
\end{itemize}
\end{frame}

\begin{frame}{Impressão}
\protect\hypertarget{impressuxe3o}{}
\begin{itemize}

\item
  \textbf{Impressão (avalia primeiro):}
  \[\frac{(\sigma, e) \to (\sigma, e')}{(\sigma,\; \mathbf{print}\; e) \to (\sigma,\; \mathbf{print}\; e')} \quad\text{(PrintStep)}\]
\end{itemize}
\end{frame}

\begin{frame}{Impressão}
\protect\hypertarget{impressuxe3o-1}{}
\begin{itemize}

\item
  \textbf{Impressão (valor final):}
  \[\frac{}{(\sigma,\; \mathbf{print}\; n) \to (\sigma,\; \mathbf{skip})} \quad\text{(Print) [emite } n\text{]}\]
\end{itemize}
\end{frame}

\begin{frame}{Sequência}
\protect\hypertarget{sequuxeancia}{}
\begin{itemize}

\item
  Skip à esquerda
\end{itemize}

\[\frac{}{(\sigma,\; \mathbf{skip};\, s^*) \to (\sigma,\; s^*)} \quad\text{(SeqSkip)}\]
\end{frame}

\begin{frame}{Sequência}
\protect\hypertarget{sequuxeancia-1}{}
\begin{itemize}

\item
  Caso geral:
\end{itemize}

\[\frac{(\sigma, s) \to (\sigma', s')}{(\sigma,\; s;\, s_1) \to (\sigma',\; s';\,s_1)} \quad\text{(Seq)}\]
\end{frame}

\begin{frame}[fragile]{Exemplo}
\protect\hypertarget{exemplo}{}
\begin{itemize}

\item
  Small-Step: \texttt{x := 2 + 3; print x * 4}

  \begin{itemize}

  \item
    Partindo do ambiente vazio \(\emptyset\):
  \end{itemize}
\end{itemize}

\[(\emptyset,\; x \mathbin{:=} 2+3;\; \mathbf{print}\; x*4)\]
\end{frame}

\begin{frame}[fragile]{Exemplo}
\protect\hypertarget{exemplo-1}{}
\begin{itemize}

\item
  Small-Step: \texttt{x := 2 + 3; print x * 4}

  \begin{itemize}

  \item
    Partindo do ambiente vazio \(\emptyset\):
  \end{itemize}
\end{itemize}

\[\xrightarrow{\text{Seq, Add}} (\emptyset,\; x \mathbin{:=} 5;\; \mathbf{print}\; x*4)\]

\[\xrightarrow{\text{Assign}} ([x \mapsto 5],\; \mathbf{skip};\; \mathbf{print}\; x*4)\]
\end{frame}

\begin{frame}[fragile]{Exemplo}
\protect\hypertarget{exemplo-2}{}
\begin{itemize}

\item
  Small-Step: \texttt{x := 2 + 3; print x * 4}

  \begin{itemize}

  \item
    Partindo do ambiente vazio \(\emptyset\):
  \end{itemize}
\end{itemize}

\[\xrightarrow{\text{SeqSkip}} ([x \mapsto 5],\; \mathbf{print}\; x*4)\]

\[\xrightarrow{\text{PrintStep, Var}} ([x \mapsto 5],\; \mathbf{print}\; 5*4)\]
\end{frame}

\begin{frame}[fragile]{Exemplo}
\protect\hypertarget{exemplo-3}{}
\begin{itemize}

\item
  Small-Step: \texttt{x := 2 + 3; print x * 4}

  \begin{itemize}

  \item
    Partindo do ambiente vazio \(\emptyset\):
  \end{itemize}
\end{itemize}

\[\xrightarrow{\text{PrintStep, Mul}} ([x \mapsto 5],\; \mathbf{print}\; 20)\]

\[\xrightarrow{\text{Print}} ([x \mapsto 5],\; \mathbf{skip}) \quad\text{[emite 20]}\]
\end{frame}

\hypertarget{semuxe2ntica-big-step}{%
\section{Semântica Big-Step}\label{semuxe2ntica-big-step}}

\begin{frame}{Expressões}
\protect\hypertarget{expressuxf5es}{}
\begin{itemize}

\item
  Nova regra: variáveis.
\end{itemize}

\[\frac{x \in \mathrm{dom}(\sigma)}{\sigma \vdash x \Downarrow \sigma(x)} \quad\text{(Var)}\]
\end{frame}

\begin{frame}{Comandos}
\protect\hypertarget{comandos}{}
\begin{itemize}

\item
  A relação \(\sigma, s \Downarrow \sigma'\) significa ``executar \(s\)
  no ambiente \(\sigma\) produz \(\sigma'\)''.
\end{itemize}

\[\frac{\sigma \vdash e \Downarrow n}
     {\sigma,\; x \mathbin{:=} e \;\Downarrow\; \sigma[x \mapsto n]}
\quad\text{(Assign)}
\]
\end{frame}

\begin{frame}{Comandos}
\protect\hypertarget{comandos-1}{}
\[
\frac{\sigma \vdash e \Downarrow n}
     {\sigma,\; \mathbf{print}\; e \;\Downarrow\; \sigma}
\quad\text{(Print) [emite } n\text{]}\]
\end{frame}

\begin{frame}{Comandos}
\protect\hypertarget{comandos-2}{}
\[\frac{}{\sigma,\; \mathbf{read}\; x \;\Downarrow\; \sigma[x \mapsto n]}
\quad\text{(Read) [} n\text{ lido da entrada]}\]
\end{frame}

\begin{frame}{Comandos}
\protect\hypertarget{comandos-3}{}
\[\frac{}{\sigma,\; \epsilon \;\Downarrow\; \sigma}
\quad\text{(SeqNil)}
\]
\end{frame}

\begin{frame}{Comandos}
\protect\hypertarget{comandos-4}{}
\[
\frac{\sigma,\; s \;\Downarrow\; \sigma' \quad \sigma',\; s^* \;\Downarrow\; \sigma''}
     {\sigma,\; s;\, s^* \;\Downarrow\; \sigma''}
\quad\text{(SeqCons)}\]
\end{frame}

\begin{frame}{Derivação: \texttt{x := 2 + 3}}
\protect\hypertarget{derivauxe7uxe3o-x-2-3}{}
\[\dfrac{
  \dfrac{
    \dfrac{}{\emptyset \vdash 2 \Downarrow 2}
    \quad
    \dfrac{}{\emptyset \vdash 3 \Downarrow 3}
  }{\emptyset \vdash 2 + 3 \Downarrow 5}
}{\emptyset,\; x \mathbin{:=} 2 + 3 \;\Downarrow\; [x \mapsto 5]}\]
\end{frame}

\begin{frame}{Derivação: \texttt{x := 2+3; print x*4}}
\protect\hypertarget{derivauxe7uxe3o-x-23-print-x4}{}
\begin{itemize}

\item
  Resolução na lousa.
\end{itemize}
\end{frame}

\hypertarget{implementauxe7uxe3o-em-haskell}{%
\section{Implementação em
Haskell}\label{implementauxe7uxe3o-em-haskell}}

\begin{frame}[fragile]{Mônada}
\protect\hypertarget{muxf4nada}{}
\begin{itemize}

\item
  A semântica big-step dita que cada comando recebe e produz um
  ambiente.
\end{itemize}

\begin{lstlisting}[language=Haskell]
type Env    = Map String Int
type Interp a = ExceptT String (StateT Env IO) a
\end{lstlisting}
\end{frame}

\begin{frame}[fragile]{Mônada}
\protect\hypertarget{muxf4nada-1}{}
\begin{itemize}

\item
  A pilha de mônadas combina:

  \begin{itemize}

  \item
    \texttt{StateT Env}: o ambiente \(\sigma\)
    implícito --- lido com \texttt{get}, modificado com
    \texttt{modify}
  \item
    \texttt{ExceptT String}: tratamento de erros
    (variável não definida)
  \item
    \texttt{IO}: efeitos de I/O para
    \texttt{read} e \texttt{print}
  \end{itemize}
\end{itemize}
\end{frame}

\begin{frame}[fragile]{Expressões}
\protect\hypertarget{expressuxf5es-1}{}
\begin{lstlisting}[language=Haskell]
interpExp :: Exp -> Interp Int
interpExp (EInt n)    = pure n
interpExp (EVar v)    = askVar v
interpExp (e1 :+: e2) = (+) <$> interpExp e1 <*> interpExp e2
interpExp (e1 :*: e2) = (*) <$> interpExp e1 <*> interpExp e2
\end{lstlisting}
\end{frame}

\begin{frame}[fragile]{Variáveis}
\protect\hypertarget{variuxe1veis}{}
\begin{lstlisting}[language=Haskell]
askVar :: String -> Interp Int
askVar s = do
    env <- get
    case Map.lookup s env of
      Nothing  -> throwError $ unwords ["Undefined variable:", s]
      Just val -> pure val
\end{lstlisting}
\end{frame}

\begin{frame}[fragile]{Comandos}
\protect\hypertarget{comandos-5}{}
\begin{lstlisting}[language=Haskell]
interpStmt :: Stmt -> Interp ()
interpStmt (SAssign v e) = do
    val <- interpExp e
    modify (Map.insert v val)
\end{lstlisting}
\end{frame}

\begin{frame}[fragile]{Comandos}
\protect\hypertarget{comandos-6}{}
\begin{lstlisting}[language=Haskell]
interpStmt (SPrint e) = do
    val <- interpExp e
    liftIO $ print val
\end{lstlisting}
\end{frame}

\begin{frame}[fragile]{Comandos}
\protect\hypertarget{comandos-7}{}
\begin{lstlisting}[language=Haskell]
interpStmt (SRead v) = do
    str <- liftIO getLine
    modify (Map.insert v (read str))

interpProgram :: Line -> Interp ()
interpProgram (Line ss) = mapM_ interpStmt ss
\end{lstlisting}
\end{frame}

\hypertarget{conclusuxe3o}{%
\section{Conclusão}\label{conclusuxe3o}}

\begin{frame}{Conclusão}
\protect\hypertarget{conclusuxe3o-1}{}
\begin{itemize}

\item
  Apresentamos a semântica operacional de uma linguagem imperativa
  simples.

  \begin{itemize}

  \item
    Foco: representação do estado.
  \item
    Small e big step
  \end{itemize}
\item
  Apresentamos a implementação em Haskell
\end{itemize}
\end{frame}
